\begin{thebibliography}{15}
\bibitem{ref1} Emmanuel Abbe, Amir Shpilka, and Avi Wigderson. Reed-Muller codes for random erasures and errors. IEEE Trans. Inform. Theory, 61(10):5229–5252, 2015. (Preliminary version in 47th STOC, 2015). doi:10.1109/TIT.2015.2462817.
\bibitem{ref2} Ido Ben-Eliezer, Rani Hod, and Shachar Lovett. Random low-degree polynomials are hard to approximate. Comput. Complexity, 21(1):63–81, 2012. (Preliminary version in 13th RANDOM, 2009). eccc:2008/TR08-080. doi:10.1007/s00037-011-0020-6.
\bibitem{ref3} Aidao Chen, Anindya De, and Aravindan Vijayaraghavan. Learning a mixture of two subspaces over finite fields. In Vitaly Feldman, Katrina Ligett, and Sivan Sabato, editors, Algorithmic Learning Theory (ALT), volume 132 of Proceedings of Machine Learning Research, pages 481–504, 2021. arXiv:2010.02841.
\bibitem{ref4} Tali Kaufman, Shachar Lovett, and Ely Porat. Weight distribution and list-decoding size of Reed-Muller codes. IEEE Trans. Inform. Theory, 58(5):2689–2696, 2012. (Preliminary version in 1st ICS, 2010). doi:10.1109/TIT.2012.2184841.
\bibitem{ref5} Peter Keevash and Benny Sudakov. Set systems with restricted cross-intersections and the minimum rank of inclusion matrices. SIAM J. Discrete Math., 18(4):713–727, 2005. doi:10.1137/S0895480103434634.
\bibitem{ref6} Shrinivas Kudekar, Santhosh Kumar, Marco Mondelli, Henry D. Pfister, Eren Eren Şaşoğlu, and Rüdiger L. Urbanke. Reed-Muller codes achieve capacity on erasure channels. IEEE Trans. Inform. Theory, 63(7):4298–4316, 2017. (Preliminary version in 48th STOC, 2016). doi:10.1109/TIT.2017.2673829.
\bibitem{ref7} David E. Muller. Application of Boolean algebra to switching circuit design and to error detection. Trans. IRE Prof. Group Electron. Comput., 3(3):6–12, 1954. doi:10.1109/IREPGELC. 1954.6499441.
\bibitem{ref8} Zipei Nie and Anthony Y. Wang. Hilbert functions and the finite degree Zariski closure in finite field combinatorial geometry. J. Comb. Theory, Ser. A, 134:196–220, 2015. doi: 10.1016/j.jcta.2015.03.011.
\bibitem{ref9} Irving S. Reed. A class of multiple-error-correcting codes and the decoding scheme. Trans. IRE Prof. Group Inf. Theory, 4:38–49, 1954. doi:10.1109/TIT.1954.1057465.
\bibitem{ref10} Galen Reeves and Henry D. Pfister. Reed-Muller codes achieve capacity on BMS channels. (manuscript), 2021. arXiv:2110.14631.
\bibitem{ref11} Alex Samorodnitsky. An upper bound on $\ell$q norms of noisy functions. IEEE Trans. Inform. Theory, 66(2):742–748, 2020. doi:10.1109/TIT.2019.2944698.
\bibitem{ref12} Ramprasad Saptharishi, Amir Shpilka, and Ben Lee Volk. Efficiently decoding Reed-Muller codes from random errors. IEEE Trans. Inform. Theory, 63(4):1954–1960, 2017. (Preliminary version in 48th STOC, 2016). doi:10.1109/TIT.2017.2671410.
\bibitem{ref13} Ori Sberlo and Amir Shpilka. On the performance of Reed-Muller codes with respect to random errors and erasures. In Shuchi Chawla, editor, Proc. 31st Annual ACM-SIAM Symp. on Discrete Algorithms (SODA), pages 1357–1376, 2020. doi:10.1137/1.9781611975994.82.
\bibitem{ref14} Claude Elwood Shannon. A mathematical theory of communication. The Bell System Technical Journal, 27:379–423, 623–656, 1948. doi:10.1002/j.1538-7305.1948.tb01338.x.
\bibitem{ref15} Victor K.-W. Wei. Generalized Hamming weights for linear codes. IEEE Trans. Inform. Theory, 37(5):1412–1418, 1991. doi:10.1109/18.133259.
\end{thebibliography}
